% !TeX root = ../TFG_LaTeX.tex
% =========================
% ===== CONFIGURACIÓN =====
% =========================

\usepackage[utf8]{inputenc}
\usepackage[catalan]{babel}

% Paquetes de mates
\usepackage{mathtools, amssymb, amsthm}

%Creació de entorns matemàtics (si vull una vegada estiga acabat arreglar numeració)
\theoremstyle{definition}
\newtheorem{definicio}{Definició}[subsection] % Definim un nou contorn de Definició, que reseteja la seua numeració cada vegada que apareix una funció.
\newtheorem*{exemple}{Exemple} % Definim un nou contorn de Exemple, que no gasta numeració (*).
\newtheorem*{exemples}{Exemples} % Definim un nou contorn de Exemples, que no gasta numeració (*).
\newtheorem{nota}[definicio]{Nota} %  Definim un nou contorn de Nota, que no gasta numeració (*).

\theoremstyle{definition}
\newtheorem{teorema}[definicio]{Teorema} % Definim un nou contorn de Teorema, que gasta la mateixa numeració que Definició.
\theoremstyle{definition}
\newtheorem{corolari}[definicio]{Corol·lari} % Definim un nou contorn de Corol·lari, que gasta la mateixa numeració que Definició.
\theoremstyle{definition}
\newtheorem{proposicio}[definicio]{Proposició} % Definim un nou contorn de Teorema, que gasta la mateixa numeració que Definició.
\newtheorem{lema}[definicio]{Lema}

\theoremstyle{remark}

%Per a que les equacions vagen numerades amb les seccions
\numberwithin{equation}{section}

% Paquets relacionats amb imatges.
\usepackage{graphicx}
\usepackage{subcaption}
\usepackage{epstopdf}

% Paquet per a personalitzar els títols i les seccions 
\usepackage{titlesec}

% Paquet per a tachar coses
\usepackage{cancel} 

% Profundidad de numeración y de índice
\setcounter{secnumdepth}{4}  % hasta \paragraph
\setcounter{tocdepth}{4}     % hasta \paragraph en el índice

% --- Numeración jerárquica ---
% (por si tu clase no lo hace ya así por defecto)
\renewcommand\theparagraph{\thesubsubsection.\arabic{paragraph}}
% Si también usas subparagraph:
% \renewcommand\thesubparagraph{\theparagraph.\arabic{subparagraph}}

% Paquets per a poder quadrar bé el text
\usepackage{float}
\usepackage{parskip}

% Paquets per a personalitzar el document.
\usepackage{xcolor}
\usepackage{enumerate}
\usepackage[left=2.54cm, right=2.54cm, top=2.54cm, bottom=2.54cm]{geometry}

% Configuració per a els encapçalaments
\usepackage{fancyhdr}
\pagestyle{fancy}

\fancyhf{}

\fancyhead[L]{\nouppercase{\rightmark}} % sección actual
\fancyfoot[C]{\thepage} % número de página

\setlength{\headheight}{16pt}

\renewcommand{\headrulewidth}{0.4pt}
\renewcommand{\sectionmark}[1]{\markright{\textsc{#1}}}


\usepackage{changepage} % Para el adjustwidth

% Secciones sin sangría extra
\titleformat{\section}
  {\normalfont\Large\bfseries}
  {\thesection}{1em}{}

\titleformat{\subsection}
  {\normalfont\large\bfseries}
  {\thesubsection}{1em}{}

  % --- Formato visual de paragraph como “mini sección” ---
\titleformat{\paragraph}
  {\normalfont\normalsize\bfseries} % estilo
  {\theparagraph}                   % muestra el número (1.1.1.1)
  {1em}                             % espacio entre número y título
  {}                                % código antes del título
\titlespacing*{\paragraph}
  {0pt}{1.5ex plus .2ex minus .2ex}{1ex} % espaciado vertical

\titlespacing*{\subsection}
{0pt}        % sin sangría
{3.5ex}      % espacio antes
{4ex}        % espacio después

% Per a que encete pàgina en cada secció
\let\oldsection\section
\renewcommand{\section}{\newpage\oldsection}

\usepackage{listings} %per a poder introduir codi

\definecolor{codegreen}{rgb}{0,0.6,0}
\definecolor{codegray}{rgb}{0.5,0.5,0.5}
\definecolor{codepurple}{rgb}{0.58,0,0.82}
\definecolor{backcolour}{rgb}{1,1,1}
\definecolor{codeblue}{rgb}{0,0.2,1}

\lstdefinestyle{mystyle}{
    backgroundcolor=\color{backcolour},   
    commentstyle=\color{codegreen},
    keywordstyle=\color{codeblue},
    numberstyle=\tiny\color{codegray},
    stringstyle=\color{codepurple},
    basicstyle=\ttfamily\footnotesize,
    breakatwhitespace=false,         
    breaklines=true,                 
    captionpos=b,                    
    keepspaces=true,                 
    numbers=left,                    
    numbersep=5pt,                  
    showspaces=false,                
    showstringspaces=false,
    showtabs=false,                  
    tabsize=2
}

\lstset{style=mystyle} %per a determinar quin estil vaig a utilitzar

% Paquet per a fer referències.

\usepackage[colorlinks=true,linkcolor=blue,urlcolor=blue,citecolor=blue]{hyperref}
%\hypersetup{colorlinks = true, linkcolor = [RGB]{100, 200, 32}} % Per a que les referències tinguen un determinat color

\usepackage{xurl}

\usepackage{caption}

% ===== FIN CONFIGURACIÓN =====